% !PRINTED 43
\weberpage{43}

%<<CONT
the exponents $\alpha$ are the same in both; then the two groups
are related to each other isomorphically.

If therefore two groups have the same invariants, we can
make these invariants the degrees of the basis elements in both
groups, and then obtain this case.

If conversely two groups are isomorphic, then the elements
of the other group corresponding to the basis elements of the
one form a basis of the latter, and hence the invariants
must also be the same in both\textsuperscript{1}).

\begin{center}
§. 11.

\gesperrt{Group Characters}.
\end{center}

It is a principal problem in the theory of \gesperrt{Abel}ian
groups to find all divisors of an \gesperrt{Abel}ian group. The
solution of this problem is essentially facilitated by the
introduction of the concept of group characters, which is also
of importance in various other respects.

Let $S$ be an \gesperrt{Abel}ian group of degree $n$, whose elements
shall be denoted by $A$, and let
\weq{1}{A_1, A_2 \ldots A_\nu}
be the elements of a basis of $S$ of degrees $a_1, a_2 \ldots a_\nu$.
Every element $A$ is therefore, and indeed in only one way,
representable in the form
\weq{2}{A = A_1^{\alpha_1} A_2^{\alpha_2} \ldots A_\nu^{\alpha_\nu}}

\textsuperscript{1}) On the theory of \gesperrt{Abel}ian groups compare:

\gesperrt{Gauss}, D\'emonstration de quelques th\'eor\`emes concernant les
périodes des classes des formes binaires du second degré. Werke,
Bd. II, S. 266.

\gesperrt{Schering}, Die Fundamentalclassen der zusammensetzbaren
arithmetischen Formen, Göttinger Abhandlungen, Bd. 14.

\gesperrt{Frobenius} and \gesperrt{Stickelberger}, Ueber Gruppen vertauschbarer
 Elemente. Crelle's Journal, Bd. 86, S. 217.

\gesperrt{Weber}, ``Theorie der \gesperrt{Abel}'schen Zahlkörper'', Acta mathematica,
Bd. 8 u. 9.

The concept of the invariants was first introduced in the memoir
by \gesperrt{Frobenius} and \gesperrt{Stickelberger}, but defined somewhat differently than
here. This older definition, which proves less suitable in its application to cyclotomic
numbers, was followed by the author in the
memoir cited in the Acta mathematica.
%CONT>>
